\chapter{The contraction theorem: an unfamiliar proof made readable}\label{ch:contraction}
\goal{Prove the contraction theorem: on a nonempty complete metric space, a map from the space to itself that multiplies every distance by at most a fixed factor below one has exactly one fixed point, and repeating the map from any starting point approaches it. Treat existence, uniqueness, convergence and the error bounds as four separate claims, each with its own short argument built from the tools of Chapter~\ref{ch:limits}. Use the error bounds to guarantee the accuracy of an approximation without knowing the exact answer, and see, by removing the hypotheses one at a time, why each of them is needed.}

\section{From solving an equation to repeating a map}\label{ch12:sec:iteration}
Section~\ref{ch04:sec:unique} proved that the equation
\[
x=\frac{x+1}{2}
\]
has exactly one real solution, namely $x=1$. Section~\ref{ch05:sec:recurrence} reached the same number in a different way. Starting from $x_0=0$, it applied the rule $x\mapsto(x+1)/2$ again and again and obtained
\[
0,\ \tfrac12,\ \tfrac34,\ \tfrac78,\ \ldots,\qquad\text{in general }x_n=1-2^{-n},
\]
and Section~\ref{ch11:sec:convergence} proved that these numbers converge to $1$. So the number $1$ can be seen in two ways: as the solution of an equation, and as the limit of a repeated update. This chapter turns the second point of view into a general method for solving equations.

The general setting is a map $T:X\to X$, where $X$ is a metric space as in Chapter~\ref{ch:limits}. (Recall from Section~\ref{ch03:sec:functions} that a map is the same thing as a function.) A \emph{fixed point} of $T$ is a point $x_*\in X$ with $T(x_*)=x_*$, that is, a point that the map leaves where it is. Section~\ref{ch05:sec:recurrence} already used this name for the number $1$ and the rule $x\mapsto(x+1)/2$. Starting from any point $x_0\in X$, we can apply $T$ over and over:
\[
x_1=T(x_0),\qquad x_2=T(x_1),\qquad x_3=T(x_2),\qquad\ldots,
\]
in general $x_{n+1}=T(x_n)$ for every $n\in\NN$. The points $x_0,x_1,x_2,\ldots$ are called the \emph{iterates} of $T$ starting from $x_0$. If the iterates converge and their limit is a fixed point, then repeating the map is a way of solving the equation $T(x)=x$. In practice we stop after finitely many steps, so what we obtain is an approximate solution, and we will want to know how good it is.

Repeating a map is a very simple algorithm, but it does not always work. Here are four ways in which it can fail.
\begin{itemize}
\item The iterates can run away. For $T(x)=2x$ on $\RR$, starting from $x_0=1$, the iterates are $1,2,4,8,\ldots$, although $T$ has the fixed point $0$.
\item The iterates can jump back and forth forever. For $T(x)=1-x$ on $\RR$, starting from $x_0=0$, the iterates are $0,1,0,1,\ldots$. The fixed point is $1/2$, since $1-\tfrac12=\tfrac12$, but every iterate stays at distance $1/2$ from it.
\item There can be several fixed points, and then the phrase ``the solution'' has no clear meaning, as Section~\ref{ch04:sec:unique} warned. For $T(x)=x^2$ on the interval $[0,1]$, both $0$ and $1$ are fixed points.
\item The iterates can approach a point that is missing from the space. Take $T(x)=x/2$ on the interval $(0,1)$. An interval like this one, which leaves out both of its endpoints, is called an \emph{open interval}. Starting from $x_0=1/2$, the iterates are $\tfrac12,\tfrac14,\tfrac18,\ldots$. They approach $0$, which does not belong to $(0,1)$, and in fact $T$ has no fixed point in $(0,1)$ at all.
\end{itemize}
We want conditions on $T$ and $X$ that rule out all of these failures in advance, so that we do not have to compute a few iterates and hope that the pattern continues. The first condition says that $T$ brings points closer together. In stating it, and from now on, we often write $Tx$ for $T(x)$ to save brackets; both mean the same point.

\par\noindent\begin{minipage}{\linewidth}
\begin{definition}[Contraction]\label{ch12:def:contraction}
Let $(X,d)$ be a metric space, and let $q$ be a real number with $0\le q<1$. A map $T:X\to X$ is a \emph{contraction with factor $q$} if
\[
d(Tx,Ty)\le q\,d(x,y)\qquad\text{for all }x,y\in X.
\]
This inequality is called the \emph{contraction inequality}, and $q$ is called a \emph{contraction factor} of $T$.
\end{definition}
\end{minipage}\par

In words: for any two points $x$ and $y$, the new points $Tx$ and $Ty$ are at most $q$ times as far apart as $x$ and $y$ were. With $q=1/2$, for instance, $Tx$ and $Ty$ are at most half as far apart as $x$ and $y$. If $T$ is a contraction with factor $q$, it is also one with every larger factor $q'$ below one: if $q\le q'<1$, then $q\,d(x,y)\le q'\,d(x,y)$, because distances are never negative. That is why we say ``a'' contraction factor rather than ``the'' contraction factor.

\subsection*{Reading the contraction inequality}
Two features of the definition matter in everything that follows.

First, $x$ and $y$ are \emph{any} two points of $X$. They need not be consecutive iterates, and neither of them needs to be close to a fixed point. For $T(x)=(x+1)/2$ on $\RR$, with the usual distance $d(x,y)=|x-y|$, we compute
\[
|T(x)-T(y)|=\left|\frac{x+1}{2}-\frac{y+1}{2}\right|
=\left|\frac{(x+1)-(y+1)}{2}\right|
=\frac12|x-y|.
\]
So this $T$ is a contraction with factor $1/2$. The constant terms cancel when two outputs are compared, just as they did in the uniqueness proof of Section~\ref{ch04:sec:unique}.

Second, one number $q$ must work for all pairs at once. It is not enough that every pair of different points moves closer together. If the amount of shrinking may depend on the pair, then the ratio $d(Tx,Ty)/d(x,y)$, for different points $x$ and $y$, can come closer and closer to $1$ as the pair changes, and then the conclusions of this chapter can fail. Section~\ref{ch12:sec:hypotheses} gives a map that brings every pair of different points strictly closer together and still has no fixed point.

\subsection*{The outputs must be allowed inputs}
The second condition is built into the arrow of $T:X\to X$: every output $T(x)$ must again be a point of $X$. A map from a set to itself is called a \emph{self-map}. This condition is what allows us to apply $T$ to $x_1$, then to $x_2$, and so on, exactly as Section~\ref{ch03:sec:functions} required the outputs of one function to be allowed inputs of the next before composing them.

The condition is easy to overlook when a formula makes sense for more numbers than we intend to use. Take $T(x)=(x+1)/2$ again, but now only on the interval $[0,1/2]$. The contraction inequality still holds for points of this interval, with factor $1/2$, because it holds for all real numbers. Yet $T(1/2)=3/4$ lies outside $[0,1/2]$, so $T$ is not a map from $[0,1/2]$ to itself. Starting from $x_0=0$, we get $x_1=1/2$ and then $x_2=3/4$, and the iterates have already left the interval. On $[0,1]$, by contrast, $T$ is a self-map: if $0\le x\le1$, then $1\le x+1\le2$, so $1/2\le T(x)\le1$.

So before using a contraction, we check two separate things: that $T$ maps $X$ into $X$, and that the contraction inequality holds for all pairs of points of $X$.

\pauseq{Is $T(x)=x^2$ a contraction of the interval $[0,1]$? Is it a contraction of the interval $[0,1/3]$?}{On $[0,1]$ it is not. It is a self-map there, because $0\le x\le1$ gives $0\le x^2\le1$. But for $x=1$ and $y=1/2$ we get $|T(x)-T(y)|=1-\tfrac14=\tfrac34$, which is larger than $|x-y|=\tfrac12$; no factor $q<1$, and not even $q=1$, can work for this pair. On $[0,1/3]$ it is a contraction. It is a self-map, because $0\le x\le1/3$ gives $0\le x^2\le1/9\le1/3$. For the contraction inequality, factor the difference of two squares: $|x^2-y^2|=|x-y|\,|x+y|$. For $x,y\in[0,1/3]$ we have $0\le x+y\le2/3$, so $|T(x)-T(y)|\le\tfrac23|x-y|$, and $q=2/3$ is a contraction factor.}

\section{The theorem and its dependency map}\label{ch12:sec:theorem}
The theorem below says that these two conditions rule out the failures of Section~\ref{ch12:sec:iteration}, provided the space is nonempty and complete. It is also known as Banach's fixed-point theorem, after the Polish mathematician Stefan Banach (1892--1945), who published it in 1922.
\par\noindent\begin{minipage}{\linewidth}
\begin{theorem}[Contraction theorem]\label{thm:banach}
Let $(X,d)$ be a nonempty complete metric space, and let $T:X\to X$ be a contraction with factor $q$, where $0\le q<1$. Then the following hold.
\begin{enumerate}
\item $T$ has exactly one fixed point $x_*$.
\item For every starting point $x_0\in X$, the iterates defined by $x_{n+1}=T(x_n)$ converge to $x_*$, and for every $n\in\NN$
\begin{equation}\label{eq:apriori}
d(x_n,x_*)\le\frac{q^n}{1-q}\,d(x_1,x_0).
\end{equation}
\item For every point $z\in X$, whether or not it is an iterate,
\begin{equation}\label{eq:residual}
d(z,x_*)\le\frac{d(z,Tz)}{1-q}.
\end{equation}
\end{enumerate}
\end{theorem}
\end{minipage}\par

Section~\ref{ch01:sec:algebra} agreed that $a^0=1$ for every nonzero number $a$; in this chapter we read $q^0$ as $1$ even when $q=0$, so that \eqref{eq:apriori} makes sense for $n=0$.

\subsection*{Four separate promises}
The theorem makes four promises, each needing its own argument.
\begin{itemize}
\item \emph{Existence}: some point of $X$ is left unchanged by $T$.
\item \emph{Uniqueness}: no two different points are left unchanged by $T$.
\item \emph{Convergence}: from every starting point, the iterates approach the fixed point.
\item \emph{Error bounds}: after finitely many steps, we can say how far we still are from the fixed point, even though we do not know where it is.
\end{itemize}
For a general map, the first two promises do not follow from each other, and even together they do not give the third. The examples of Section~\ref{ch12:sec:iteration} show this: $x\mapsto x^2$ on $[0,1]$ has fixed points, but more than one; $x\mapsto x/2$ on $(0,1)$ has at most one, but only because it has none; and $x\mapsto1-x$ on $\RR$ has exactly one, yet its iterates from $x_0=0$ never approach it.

\pauseq{The map $T(x)=1-x$ on $\RR$ has a fixed point, but its iterates from $x_0=0$ do not converge. Which hypothesis of Theorem~\ref{thm:banach} does it fail?}{The space $\RR$ is nonempty and complete, and $T$ maps $\RR$ to itself. The contraction inequality is what fails: $|T(x)-T(y)|=|(1-x)-(1-y)|=|x-y|$, so $T$ never shrinks a distance, and no factor $q<1$ works.}

\subsection*{Two ways to measure the error}
The \emph{error} of a point $z\in X$ is its distance $d(z,x_*)$ from the fixed point. Parts~2 and~3 of the theorem bound it in two ways.

The bound \eqref{eq:apriori} is called the \emph{a priori estimate}, from the Latin for ``from what comes before.'' It uses only $q$ and the length $d(x_1,x_0)$ of the first step, so after a single step it tells us how many steps will be enough for a desired accuracy.

The bound \eqref{eq:residual} is called the \emph{residual estimate}. The number $d(z,Tz)$ is the \emph{residual} of $z$: it measures how far $z$ is from satisfying $Tz=z$, and it is zero exactly when $z$ is a fixed point. The point $z$ may be an iterate, a guess, or the output of any other method; one evaluation of $T$ gives a guaranteed upper bound on its error. A guaranteed bound of this kind, computed from known quantities, is called an \emph{error certificate}.

Both bounds divide by $1-q$, which is positive because $q<1$; when $q$ is close to $1$, the factor $1/(1-q)$ is large and the bounds are weak.

\subsection*{The plan of the proof}
Here is the proof in outline, with the earlier facts that each step uses.
\begin{enumerate}
\item \emph{The steps shrink geometrically.} Nonemptiness gives a starting point, and the self-map condition lets us form all the iterates. The contraction inequality and induction (Chapter~\ref{ch:induction}) show that the step from $x_n$ to $x_{n+1}$ is at most $q^n$ times as long as the first step.
\item \emph{The iterates form a Cauchy sequence,} by Proposition~\ref{ch11:prop:geometric-cauchy}, which rests on the triangle inequality and the geometric sum.
\item \emph{A limit $x_*$ exists in $X$,} by completeness (Section~\ref{ch11:sec:complete}).
\item \emph{The limit is a fixed point.} The triangle inequality, the contraction inequality and the definition of convergence show that $d(Tx_*,x_*)$ is smaller than every positive tolerance, hence zero.
\item \emph{There is only one fixed point.} The contraction inequality, applied to two fixed points, shows that they coincide: the comparison argument of Section~\ref{ch04:sec:unique} in a new setting.
\item \emph{The error bounds hold.} The triangle inequality, the contraction inequality and $Tx_*=x_*$ give \eqref{eq:residual}; applied to $x_n$, together with Step~1, it gives \eqref{eq:apriori}.
\end{enumerate}
A list like this, which records what each step of a proof depends on, is a \emph{dependency map} of the proof. It tells you what you need to know before reading the proof, and where each hypothesis is used. Here everything comes from Chapters~\ref{ch:induction} and~\ref{ch:limits}; no calculus is needed. Chapter~\ref{ch:bridge} uses dependency maps to approach a proof whose ingredients are not all familiar yet.

\section{The proof, one step at a time}\label{ch12:sec:proof}
The proof follows the six steps of Section~\ref{ch12:sec:theorem}.

\begin{proof}
\textbf{Step 1: the steps shrink geometrically.} Since $X$ is nonempty, we can choose a starting point $x_0\in X$, any point at all. Define $x_{n+1}=T(x_n)$ for every $n\in\NN$; because $T$ maps $X$ into $X$, each iterate is again a point of $X$, so $T$ can be applied to it. Write
\[
D=d(x_1,x_0)
\]
for the length of the first step. We prove by induction on $n$ that
\[
d(x_{n+1},x_n)\le q^nD\qquad\text{for every }n\in\NN.
\]
For $n=0$ both sides equal $D$, because $q^0=1$. If the inequality holds for some $n$, then, since $x_{n+2}=T(x_{n+1})$ and $x_{n+1}=T(x_n)$, the contraction inequality for the points $x_{n+1}$ and $x_n$ gives
\[
d(x_{n+2},x_{n+1})=d(Tx_{n+1},Tx_n)
\le q\,d(x_{n+1},x_n)\le q\cdot q^nD=q^{n+1}D,
\]
where the last inequality multiplies the induction hypothesis by $q\ge0$. This is the claim for $n+1$, which completes the induction.

\textbf{Step 2: the iterates form a Cauchy sequence.} Step~1 is exactly the hypothesis $d(x_{j+1},x_j)\le Dq^j$ of Proposition~\ref{ch11:prop:geometric-cauchy}, with $D=d(x_1,x_0)\ge0$. The proposition assumes $0<q<1$. If $q=0$, then $T$ is also a contraction with factor $1/2$ (Section~\ref{ch12:sec:iteration}), so Step~1 with $1/2$ in place of $q$ gives $d(x_{j+1},x_j)\le D(1/2)^j$. Either way, the proposition shows that the iterates form a Cauchy sequence.

\textbf{Step 3: completeness supplies a limit.} Since $X$ is complete, the Cauchy sequence $x_0,x_1,x_2,\ldots$ converges to some point $x_*$ of $X$. This is the only place in the proof where completeness is used.

\textbf{Step 4: the limit is a fixed point.} We show that $d(Tx_*,x_*)=0$. For every $n\in\NN$, the triangle inequality with $Tx_n$ as a stop on the way, the contraction inequality for the pair $x_*,x_n$, and $Tx_n=x_{n+1}$ give
\[
d(Tx_*,x_*)\le d(Tx_*,Tx_n)+d(Tx_n,x_*)
\le q\,d(x_*,x_n)+d(x_{n+1},x_*).
\]
Both distances on the right become small when $n$ is large. To make this precise, let $\eps>0$. Because $x_n\to x_*$, the definition of convergence with tolerance $\eps/2$ gives a threshold $N$ such that $d(x_k,x_*)<\eps/2$ for every $k\ge N$. Take $n=N$ in the display; then both $n$ and $n+1$ are at least $N$. Since $q<1$, we have $q\,d(x_*,x_n)\le d(x_n,x_*)<\eps/2$, and also $d(x_{n+1},x_*)<\eps/2$. Hence
\[
d(Tx_*,x_*)<\frac\eps2+\frac\eps2=\eps.
\]
So the number $r=d(Tx_*,x_*)$, which does not depend on $n$ or on $\eps$, is smaller than every positive $\eps$. Such a number must be zero: $r\ge0$, and if $r$ were positive, taking $\eps=r$ would give $r<r$. So $d(Tx_*,x_*)=0$, and by the first requirement of a metric, $Tx_*=x_*$.

\textbf{Step 5: there is only one fixed point.} Suppose $y_*$ is also a fixed point of $T$. Since $x_*=Tx_*$ and $y_*=Ty_*$, the contraction inequality gives
\[
d(x_*,y_*)=d(Tx_*,Ty_*)\le q\,d(x_*,y_*).
\]
Subtracting $q\,d(x_*,y_*)$ from both sides gives $(1-q)\,d(x_*,y_*)\le0$, and dividing by the positive number $1-q$ gives $d(x_*,y_*)\le0$. Distances are never negative, so $d(x_*,y_*)=0$, and therefore $y_*=x_*$.

This also proves convergence from \emph{every} starting point: $x_0$ in Step~1 was arbitrary, Steps~1--4 showed that its iterates converge to a fixed point of $T$, and by Step~5 that fixed point can only be $x_*$.

\textbf{Step 6: the error bounds.} We prove the residual estimate first. Let $z$ be any point of $X$. Use the triangle inequality with $Tz$ as a stop, then $x_*=Tx_*$, then the contraction inequality:
\[
d(z,x_*)\le d(z,Tz)+d(Tz,x_*)=d(z,Tz)+d(Tz,Tx_*)
\le d(z,Tz)+q\,d(z,x_*).
\]
Subtracting $q\,d(z,x_*)$ gives $(1-q)\,d(z,x_*)\le d(z,Tz)$, and dividing by $1-q>0$ gives \eqref{eq:residual}.

For the a priori estimate, fix $n$ and take $z=x_n$. Its residual is $d(x_n,Tx_n)=d(x_n,x_{n+1})=d(x_{n+1},x_n)$, which Step~1 bounds by $q^nD$. So the residual estimate gives
\[
d(x_n,x_*)\le\frac{d(x_n,x_{n+1})}{1-q}\le\frac{q^nD}{1-q}=\frac{q^n}{1-q}\,d(x_1,x_0),
\]
which is \eqref{eq:apriori}. This completes the proof.
\end{proof}

\pauseq{Use the residual estimate \eqref{eq:residual} to give a second, one-line proof that $T$ has at most one fixed point.}{If $z$ is any fixed point, then $Tz=z$, so its residual $d(z,Tz)$ is $0$, and the residual estimate gives $d(z,x_*)\le0/(1-q)=0$. So $z=x_*$.}

\section{What the error estimates are for}\label{ch12:sec:estimates}
\subsection*{Predicting how many steps are enough}
Take $T(x)=(x+1)/2$ on $\RR$, a contraction with factor $q=1/2$, and start from $x_0=0$. Then $x_1=1/2$, so $D=d(x_1,x_0)=1/2$, and the a priori estimate \eqref{eq:apriori} becomes
\[
|x_n-1|\le\frac{(1/2)^n}{1-1/2}\cdot\frac12=2^{-n}.
\]
Suppose we want an error below $0.001=1/1000$. Since $2^9=512$ and $2^{10}=1024$, the bound $2^{-n}$ is still larger than $1/1000$ at $n=9$ but equals $1/1024$ at $n=10$. So ten steps are guaranteed to be enough. Here the estimate cannot be improved: by the formula $x_n=1-2^{-n}$, the error is exactly $2^{-n}$, so nine steps really are too few. Usually, though, the true error is smaller than this worst-case bound, so the estimate may ask for more steps than necessary.

\subsection*{Certifying an approximation found by any method}
To certify a point $z$, the residual estimate \eqref{eq:residual} multiplies its residual by $1/(1-q)$, which can be large; the residual by itself is not a bound on the error. For example, $T(x)=0.99x+0.01$ is a contraction of $\RR$ with factor $q=0.99$, since $|T(x)-T(y)|=0.99\,|x-y|$, and its fixed point is $1$. The point $z=0$ has the small residual $|0-T(0)|=0.01$, yet its error is $|0-1|=1$, a hundred times larger. The residual estimate gives the bound $0.01/(1-0.99)=1$, exactly the true error. So a small residual guarantees a small error only when $q$ is not too close to $1$.

\subsection*{Using the length of the latest step}
For an iterate $x_n$ with $n\ge1$, there is a third bound, which uses the step just taken:
\[
d(x_n,x_*)\le\frac{q}{1-q}\,d(x_n,x_{n-1}).
\]
To prove it, apply the residual estimate to $z=x_n$. Its residual is $d(x_n,Tx_n)=d(x_n,x_{n+1})$, and since $x_n=Tx_{n-1}$ and $x_{n+1}=Tx_n$, the contraction inequality gives
\[
d(x_n,x_{n+1})=d(Tx_{n-1},Tx_n)\le q\,d(x_{n-1},x_n)=q\,d(x_n,x_{n-1}).
\]
So $d(x_n,x_*)\le d(x_n,x_{n+1})/(1-q)\le q\,d(x_n,x_{n-1})/(1-q)$.

This is called an \emph{a posteriori estimate}, from the Latin for ``from what comes after,'' because it can be evaluated only after the step from $x_{n-1}$ to $x_n$ has been computed. It is never worse than the a priori estimate, because Step~1 of the proof gives $d(x_n,x_{n-1})\le q^{n-1}D$, so
\[
\frac{q}{1-q}\,d(x_n,x_{n-1})\le\frac{q\cdot q^{n-1}D}{1-q}=\frac{q^n}{1-q}\,D,
\]
which is the a priori bound. When the observed steps shrink faster than the worst case $q^{n-1}D$, the a posteriori estimate is much better.

\pauseq{For $T(x)=(x+1)/2$ and $x_0=0$, compute the a posteriori bound for $x_3$ and compare it with the true error $|x_3-1|$.}{The iterates are $x_2=3/4$ and $x_3=7/8$, so the latest step has length $1/8$. With $q=1/2$ we have $q/(1-q)=1$, so the bound is $|x_3-1|\le1/8$. The true error is $1-7/8=1/8$, so the bound is attained, as the a priori bound $2^{-3}$ is for this map.}

\section{A quadratic example, worked by hand}\label{ch12:sec:example}
For $T(x)=(x+1)/2$ we knew the fixed point from the start. Now we do not, and the theorem tells us about it before we have found it. Let
\[
T(x)=\frac{x^2+1}{3}\qquad\text{on }X=[0,1],
\]
with the usual distance $|x-y|$. First we check each hypothesis.
\begin{itemize}
\item \emph{The space.} The interval $[0,1]$ is nonempty, and complete by Section~\ref{ch11:sec:complete}.
\item \emph{The self-map condition.} If $0\le x\le1$, then $0\le x^2\le1$, so $1\le x^2+1\le2$ and $1/3\le T(x)\le2/3$. In particular $T(x)\in[0,1]$.
\item \emph{The contraction inequality.} For $x,y\in[0,1]$, factor the difference of two squares and use $|ab|=|a|\,|b|$ and $0\le x+y\le2$:
\[
|T(x)-T(y)|=\frac{|x^2-y^2|}{3}=\frac{|x-y|\,|x+y|}{3}\le\frac23|x-y|.
\]
So $q=2/3$ is a contraction factor.
\end{itemize}
By Theorem~\ref{thm:banach}, $T$ has exactly one fixed point $x_*$ in $[0,1]$, and the iterates from every starting point in $[0,1]$ converge to it.

\subsection*{Iterates and a certificate}
Start from $x_0=0$, and compute with exact fractions, not rounded decimals:
\begin{align*}
x_1&=\frac{0^2+1}{3}=\frac13,\\
x_2&=\frac{(1/3)^2+1}{3}=\frac{1/9+9/9}{3}=\frac{10/9}{3}=\frac{10}{27},\\
x_3&=\frac{(10/27)^2+1}{3}=\frac{100/729+729/729}{3}=\frac{829/729}{3}=\frac{829}{2187}.
\end{align*}
Take $z=x_2$. Then $Tz=x_3$, so the residual of $z$ is $|x_2-x_3|$. Writing $x_2=810/2187$ over the denominator $2187=27\cdot81$,
\[
d(z,Tz)=\left|\frac{810}{2187}-\frac{829}{2187}\right|=\frac{19}{2187}.
\]
Since $1/(1-q)=1/(1-2/3)=3$, the residual estimate \eqref{eq:residual} gives the certificate
\[
d(x_2,x_*)\le3\cdot\frac{19}{2187}=\frac{19}{729}=0.02606\ldots.
\]
A certificate rounded to a decimal must be rounded up, never down; otherwise it is no longer a guarantee. Rounded up, the certificate says that $x_2=10/27\approx0.370$ lies within $0.027$ of the fixed point, although we do not yet know where the fixed point is.

The other two estimates of Section~\ref{ch12:sec:estimates} give weaker bounds. The a priori estimate, with $D=d(x_1,x_0)=1/3$ and $n=2$, uses only the first step:
\[
d(x_2,x_*)\le\frac{(2/3)^2}{1-2/3}\cdot\frac13=\frac49=0.444\ldots.
\]
The a posteriori estimate uses the step from $x_1=9/27$ to $x_2=10/27$ and gives $d(x_2,x_*)\le\frac{2/3}{1/3}\cdot\frac1{27}=\frac2{27}=0.0740\ldots$. The residual estimate is the sharpest, because it uses the extra information in $x_3$.

\subsection*{Checking against the exact solution}
Here we can also find $x_*$ exactly. The fixed-point equation $(x^2+1)/3=x$ becomes $x^2+1=3x$, that is, $x^2-3x+1=0$. Multiply by $4$ and add $5$ to both sides to complete the square:
\[
4x^2-12x+9=5,\qquad\text{that is,}\qquad(2x-3)^2=5.
\]
If $a^2=5$, then $(a-\sqrt5)(a+\sqrt5)=a^2-5=0$, so $a=\pm\sqrt5$. Hence $2x-3=\pm\sqrt5$, and $x=(3\pm\sqrt5)/2$. Since $4<5<9$, we have $2<\sqrt5<3$, so $(3+\sqrt5)/2>5/2$ lies outside $[0,1]$, while $(3-\sqrt5)/2$ lies between $0$ and $1/2$. The fixed point in $[0,1]$ is therefore
\[
x_*=\frac{3-\sqrt5}{2}\approx0.382.
\]
The true error of $x_2$ is $x_*-10/27\approx0.0116$, comfortably below the certificate.

On all of $\RR$, the same formula has the two fixed points $(3\pm\sqrt5)/2$. This does not contradict uniqueness: Step~5 shows that a contraction never has two different fixed points, so $T$ is not a contraction of $\RR$. The uniqueness in Theorem~\ref{thm:banach} is about fixed points that lie in $X$.

\section{Remove the hypotheses one at a time}\label{ch12:sec:hypotheses}
Theorem~\ref{thm:banach} has four hypotheses: $X$ is nonempty, $X$ is complete, $T$ maps $X$ into itself, and $T$ satisfies the contraction inequality with a factor $q<1$. A good way to understand why a hypothesis is there is to remove it, keep all the others, and look for an example in which the conclusion fails. Exercise~\ref{ex:12.4} asks you to check that each example in the first three groups below satisfies the hypotheses it keeps.

\subsection*{Without completeness}
Let $T(x)=x/2$ on the open interval $X=(0,1)$. This $T$ is a self-map and a contraction with factor $1/2$, but it has no fixed point in $X$: the equation $x/2=x$ forces $x=0$, which does not belong to $(0,1)$. The proof breaks down at Step~3. Starting from $x_0=1/2$, Steps~1 and~2 go through, and the iterates $x_n=2^{-(n+1)}$ form exactly the Cauchy sequence that Section~\ref{ch11:sec:complete} used to show that $(0,1)$ is not complete. Its only possible limit, $0$, is missing from the space, so there is no point $x_*$ for Steps~4--6 to work with.

\subsection*{Without a factor below one}
Suppose we allow $q=1$, so that the inequality only says $d(Tx,Ty)\le d(x,y)$. On $\RR$, the identity map $T(x)=x$ satisfies this with equality, and every real number is a fixed point, so uniqueness fails. The map $T(x)=x+1$ also satisfies $|T(x)-T(y)|=|x-y|$, but $x+1=x$ has no solution, so existence fails. So $q=1$ does not merely slow convergence down; it can destroy existence or uniqueness. In the proof, Step~2 relies on the geometric tail bound $q^n/(1-q)$, and Steps~5 and~6 divide by $1-q$, which is now zero.

Even a map that brings every pair of different points strictly closer together can fail, if the shrinking gets weaker and weaker. Let $X=[1,\infty)$, the set of real numbers $x\ge1$, and let
\[
T(x)=x+\frac1x.
\]
For $x\ge1$ we have $T(x)>x\ge1$, so $T$ maps $X$ into itself, and the set $[1,\infty)$ is complete, by an adaptation of the argument given for $[0,1]$ in Section~\ref{ch11:sec:complete}. Since $\frac1x-\frac1y=-\frac{x-y}{xy}$,
\[
T(x)-T(y)=(x-y)+\left(\frac1x-\frac1y\right)=(x-y)\left(1-\frac1{xy}\right).
\]
For $x,y\ge1$ with $x\ne y$, at least one of the two numbers is larger than $1$, so $xy>1$ and the factor $1-1/(xy)$ lies strictly between $0$ and $1$. Hence $|T(x)-T(y)|<|x-y|$. But the factor depends on the pair, and no single $q<1$ works for all pairs. Indeed, let $0\le q<1$, and use the Archimedean property (Section~\ref{ch02:sec:quantifiers}) to choose an integer $k\ge1$ with $k>1/(1-q)$. For $x=k+1$ and $y=k$ the ratio $|T(x)-T(y)|/|x-y|$ is $1-1/(k(k+1))$. Since $k(k+1)=k^2+k>k>1/(1-q)$, taking reciprocals gives $1/(k(k+1))<1-q$, so the ratio is larger than $q$, and $q$ is not a contraction factor. Finally, $T$ has no fixed point: $x+1/x=x$ would require $1/x=0$.

\subsection*{Without the self-map condition}
Let $T(x)=x/2+1$, used only for inputs in the nonempty complete space $[0,1]$. The contraction inequality holds with factor $1/2$, but $T(1)=3/2$ lies outside $[0,1]$, so $T$ is not a map from $[0,1]$ to itself. The only real solution of $x/2+1=x$ is $x=2$, so $T$ has no fixed point in $[0,1]$. The proof breaks down already in Step~1: from $x_0=0$ we get $x_1=1$ and then $x_2=3/2$, which has left the space, so $T(x_2)$ is not even defined.

\subsection*{Without nonemptiness}
If $X$ is the empty set, there is still a map $T:X\to X$: it has no inputs, so there is nothing to assign. The requirements of a metric, of completeness and of the contraction inequality all hold vacuously (Section~\ref{ch02:sec:contrapositive}), because there are no points to check them for. But there is no fixed point either. The proof uses nonemptiness only to choose $x_0$ in Step~1.

\subsection*{Summary}
The table records where the proof uses each hypothesis, and examples showing that it cannot be dropped.
\begin{center}
\begin{tabular}{@{}lll@{}}\toprule
Hypothesis & Used in the proof & Examples without it\\\midrule
$X$ is nonempty & Step~1, to choose $x_0$ & the empty set\\
$X$ is complete & Step~3 & $x\mapsto x/2$ on $(0,1)$\\
$T$ maps $X$ into $X$ & Step~1, to form the iterates & $x\mapsto x/2+1$ on $[0,1]$\\
factor $q<1$ & Steps~2, 5 and~6 & $x\mapsto x$ and $x\mapsto x+1$ on $\RR$\\
 & & $x\mapsto x+1/x$ on $[1,\infty)$\\
\bottomrule
\end{tabular}
\end{center}

\section*{Exercises}
Exercises~\ref{ex:12.1}--\ref{ex:12.4} are core exercises, and Exercises~\ref{ex:12.5} and~\ref{ex:12.6} are transfer exercises. Complete solutions are in Appendix~\ref{app:solb}, under ``Week~12.''
\begin{exercise}\label{ex:12.1}Let $T(x)=(3x+1)/4$ on $\RR$. Find the fixed point of $T$ and the smallest contraction factor of $T$. Then find a formula for the iterates $x_n$ starting from $x_0=0$, and prove it by induction.\end{exercise}
\begin{exercise}\label{ex:12.2}Let $T$ be a contraction with factor $q=3/4$ of a nonempty complete metric space, with fixed point $x_*$, and let $z$ be a point with $d(z,Tz)\le0.002$. What upper bound on $d(z,x_*)$ does Theorem~\ref{thm:banach} guarantee? Is the residual bound $0.002$ itself an upper bound on $d(z,x_*)$ in every such case? Give a reason or an example.\end{exercise}
\begin{exercise}\label{ex:12.3}Let $T(x)=(x^2+1)/3$. Show that $T$ maps the interval $[1/3,2/3]$ into itself and that $4/9$ is a contraction factor of $T$ on this interval. Starting from $x_0=0$, explain why every iterate from $x_1$ on lies in $[1/3,2/3]$. Use the smaller factor to improve the certificate $d(x_2,x_*)\le19/729$ obtained in Section~\ref{ch12:sec:example}, and explain why the improvement is justified.\end{exercise}
\begin{exercise}\label{ex:12.4}Section~\ref{ch12:sec:hypotheses} gives counterexamples in three groups: without completeness, without a factor below one, and without the self-map condition. For each example in these three groups, prove that it satisfies every hypothesis of Theorem~\ref{thm:banach} except the one named in the heading of its group. For $T(x)=x+1/x$, this includes proving that $[1,\infty)$ is complete.\end{exercise}
\begin{exercise}\label{ex:12.5}Let $a$ and $b$ be real numbers with $|a|<1$, and let $T(x)=ax+b$ on $\RR$. Find the fixed point $x_*$ of $T$, and prove by induction that for every starting point $x_0$ the iterates satisfy $x_n-x_*=a^n(x_0-x_*)$ for every $n\in\NN$. How do the iterates move when $a<0$, compared with when $0<a<1$?\end{exercise}
\begin{exercise}\label{ex:12.6}Let $T$ and $S$ be two self-maps of the same nonempty complete metric space $(X,d)$, both contractions with factor $q$, where $0\le q<1$. Write $x_T$ and $x_S$ for their fixed points. Suppose there is a number $\eta\ge0$ (the Greek letter eta) with $d(Tx,Sx)\le\eta$ for every $x\in X$. Prove that $d(x_T,x_S)\le\eta/(1-q)$.\end{exercise}

